% Part: propositional-logic
% Chapter: syntax-and-semantics
% Section: semantic-notions

\documentclass[../../../include/open-logic-section]{subfiles}

\begin{document}

\olfileid[hi]{pl}{syn}{sem}

\olsection{अर्थवैज्ञानिक धारणाएँ}

हम निम्नलिखित अर्थवैज्ञानिक धारणाएँ परिभाषित करते हैं:

\begin{defn} 
\begin{enumerate}
\item !!^a{formula}~$!A$ \emph{संतुष्टनीय} है यदि किसी
  ऐसे~$\pAssign{v}$ के लिए $\pSat{v}{!A}$ हो; वह
  \emph{असंतुष्टनीय} है यदि किसी भी $\pAssign{v}$ के लिए $\pSat{v}{!A}$ न हो;
\item !!^a{formula}~$!A$ \emph{पुनरुक्ति} है यदि $\pSat{v}{!A}$ सभी
  !!{valuation}s~$\pAssign{v}$ के लिए सत्य हो;
\item !!^a{formula}~$!A$ \emph{आपातिक} है यदि वह संतुष्टनीय है लेकिन
  पुनरुक्ति नहीं है;
\item यदि $\Gamma$ !!{formula}s का कोई समुच्चय है, तो $\Gamma \Entails !A$
  (``$\Gamma$ से $!A$ तार्किकतः निष्पन्न है'') तभी और केवल तभी, जब
  $\pSat{v}{!A}$ उन सभी !!{valuation}~$\pAssign{v}$ के लिए सत्य हो जिनके लिए
  $\pSat{v}{\Gamma}$ हो।
\item यदि $\Gamma$ !!{formula}s का कोई समुच्चय है, तो $\Gamma$
  \emph{संतुष्टनीय} है यदि कोई !!a{valuation}~$\pAssign{v}$ ऐसा हो जिसके लिए
  $\pSat{v}{\Gamma}$ हो; अन्यथा $\Gamma$
  \emph{असंतुष्टनीय} है।
\end{enumerate} 
\end{defn}

\begin{prob} 
 नीचे दिए चारों !!{formula}s में प्रत्येक के लिए बताएँ कि वह (क)~संतुष्टनीय,
 (ख)~पुनरुक्ति और (ग)~आपातिक है या नहीं।
\begin{enumerate}
  \item \( ( \Obj p_0 \lif ( \lnot \Obj p_1 \lif \lnot \Obj p_0 ) ) \).
  \item \( ( ( \Obj p_0 \land \lnot \Obj p_1 ) \lif ( \lnot \Obj p_0 \land \Obj p_2 )) \liff ( ( \Obj p_2 \lif \Obj p_0 ) \lif ( \Obj p_0 \lif \Obj p_1 )) \).
  \item \( ( \Obj p_0 \liff \Obj p_1 ) \lif ( \Obj p_2 \liff \lnot \Obj p_1 ) \).
  \item \( (( \Obj p_0 \liff ( \lnot \Obj p_1 \land \Obj p_2 )) \lor ( \Obj p_2 \lif ( \Obj p_0 \liff \Obj p_1 ))) \).
\end{enumerate}
\end{prob}

\begin{prop}
\ollabel{prop:semanticalfacts} 
\begin{enumerate} 
\item $!A$ तभी और केवल तभी पुनरुक्ति है, जब
  $\emptyset \Entails !A$; 
\item यदि $\Gamma \Entails !A$ और $\Gamma \Entails !A \lif !B$, तो
  $\Gamma \Entails !B$;
\item यदि $\Gamma$ संतुष्टनीय है, तो $\Gamma$ का प्रत्येक परिमित उपसमुच्चय
  भी संतुष्टनीय है; 
\item \ollabel{def:monotonicity} एकदिष्टता: यदि $\Gamma \subseteq \Delta$
  और $\Gamma \Entails !A$, तो $\Delta \Entails !A$ भी;
\item \ollabel{def:Cut} संक्रामकता: यदि $\Gamma \Entails !A$ और
  $\Delta \cup \{ !A\} \Entails !B$, तो $\Gamma \cup \Delta \Entails
  !B$।
\end{enumerate}
\end{prop}

\begin{proof}
अभ्यास।
\end{proof}

\begin{prob}
\olref[pl][syn][sem]{prop:semanticalfacts} सिद्ध करें।
\end{prob}

\begin{prop}\ollabel{prop:entails-unsat}
  $\Gamma \Entails !A$ तभी और केवल तभी, जब $\Gamma \cup \{\lnot !A\}$
  असंतुष्टनीय हो।
\end{prop}

\begin{proof}
अभ्यास।
\end{proof}

\begin{prob}
\olref[pl][syn][sem]{prop:entails-unsat} सिद्ध करें।
\end{prob}

\begin{thm}[अर्थवैज्ञानिक निगमन प्रमेय]
  \ollabel{thm:sem-deduction} $\Gamma \Entails !A \lif !B$ तभी और केवल
  तभी, जब $\Gamma \cup \{!A\} \Entails !B$।
\end{thm}

\begin{proof}
अभ्यास।
\end{proof}

\begin{prob}
\olref[pl][syn][sem]{thm:sem-deduction} सिद्ध करें।
\end{prob}

\end{document}
